\chapter{Engineering Scope and Evidence}

\section{Scope boundaries}
Define what is included and excluded. A strong scope statement prevents the public
report from implying fabrication, certification, field validation, or production
readiness that was not actually achieved.

\begin{table}[H]
\centering
\caption{Example scope boundary table.}
\label{tab:scope}
\begin{tabularx}{\textwidth}{>{\bfseries}p{0.22\textwidth}XX}
\toprule
Area & Included & Excluded or deferred \\
\midrule
Analytical work & Governing equations, first-order sizing, sensitivity checks & Formal uncertainty propagation \\
Numerical work & Baseline model, parameter sweep, selected convergence checks & Full certification-grade validation \\
Hardware & Interface and test concept & Fabricated prototype and chamber measurements \\
Public release & Verified plots, scripts, assumptions, limitations & Proprietary files and institution-specific administration \\
\bottomrule
\end{tabularx}
\end{table}

\section{Evidence taxonomy}
Use evidence labels consistently throughout the report:

\begin{itemize}[leftmargin=2em]
\item \AnalyticalTag\ results derived from equations, calculations, or independent scripts;
\item \SimulatedTag\ results exported from a numerical model or digital twin;
\item \MeasuredTag\ results obtained from physical instrumentation or test equipment;
\item \QualifiedTag\ claims that remain dependent on stated assumptions or incomplete validation.
\end{itemize}

\section{Deliverables}
List the public artifacts that allow a reviewer to reproduce or audit the work: source
report, code, console output, model screenshots, run manifest, parameter tables,
verification matrix, and selected native project files when licensing and size permit.
